\subsubsection{Iván Pastor Sacristán}

\begin{enumerate}
	\item Investigación del estado del arte: Investigué sobre el estado actual de la detección de deepfakes. Durante la investigación leí dos papers. En el primero hablaban sobre adaptación de los modelos a nuevos conjuntos de datos distintos de los que se usaron para entrenarlo y probarlo. En el segundo paper se centraban en un modelo capaz de clasificar correctamente los audios, aún y cuando los generados por IA habían sido modificados tanto en el dominio de la frecuencia como en el del tiempo para añadirles defectos característicos de los audios humanos como el ruido. %Añadir los enlaces a los papers en la bibliografía
	
	\item Prueba del dataset WaveFake: Hice una primera versión del extractor de características entre las que se incluyen rms, zcr, pitch, mfcc's, duración, transcripción, espectral centroid, espectral rolloff y el espectral bandwidth.
	
	\item Prueba del dataset ASVSpoof2021: Usé el mismo extractor de características e hice un clasificador difuso en el que, basándome en las características extraídas, clasifique los audios por género y edad aparente.
	
	\item Modelo compuesto: El modelo se basa en hacer un embedding de las características con un perceptrón multicapa (MLP), hacer un embedding del espectrograma del audio con una red neuronal convolucional (CNN), concatenarlos, y luego usar esos embeddings como entrada de un MLP + clasificador difuso que es lo que se encargará realmente de la tarea de clasificación.
	
	\item Fine-tuning Wav2Vev2: Usé el modelo desarrollado por Meta Wav2Vec2 y le hice un fine-tuning con nuestro dataset para que se enfocara en la clasificación de audios. Luego procesé la salida del modelo para pasar de la tupla que me devuelve el modelo a un valor entre 0 y 1 que indique la confianza de que un audio sea generado.
	
	\item XAI sobre Wav2Vec2: Sobre el modelo Wav2Vec2, no se puede aplicar directamente IA explicable ya que la entrada son embeddings, para ello lo que he hecho ha sido dividir entre las características y el espectrograma de Mel. Para las características usa Shap y Lime sobre surrogate models que imitan la salida de Wav2Vec2 y para los espectrogramas de Mel use una CNN a la que le apliqué Shap, Lime, Integrated gradients y Grad Cam.

